\section{Additional Derivations and Results}

\subsection{Proof of Theorem~\ref{thm:fixed-h-dml}}

Let $\widehat\Psi_n(\theta,\eta)$ denote the cross-fitted sample moment. Expand
the moment around the population pair $\theta_{0,h}$ and $\eta_{0,h}$ to
obtain
\begin{equation}
0=\widehat\Psi_n(\theta_{0,h},\eta_{0,h})
+J_h(\widehat\theta_h-\theta_{0,h})+R_{1n}+R_{2n},
\end{equation}
where $R_{1n}$ is the empirical-process effect of replacing $\eta_{0,h}$ by
the fold-specific nuisance estimate and $R_{2n}$ is its population remainder.
More explicitly, write $P_{n,k}$ for the empirical measure on holdout fold
$k$ and $P$ for the population measure. At $\theta_{0,h}$, the empirical-
process term is
\begin{equation}
R_{1n}=\sum_{k=1}^K\frac{|I_k|}{n}(P_{n,k}-P)
\left[\psi_h(\widehat\eta_h^{(-k)})
-\psi_h(\eta_{0,h})\right],
\end{equation}
where the common arguments $(B;\theta_{0,h})$ are suppressed. Conditional on
the training sample, cross-fitting makes each bracketed function fixed and
independent of its holdout observations. Bounded first derivatives of the
score and the nuisance-rate condition imply that its conditional second
moment is $O_p(\rho_n^2)$. With fixed $K$, conditional Chebyshev's inequality
therefore yields
\begin{equation}
R_{1n}=O_p(\rho_n/\sqrt n)=o_p(n^{-1/2}).
\end{equation}
Neyman orthogonality sets the first Gateaux derivative of the population score
to zero. A fold-by-fold second-order Taylor expansion and the bounded second
derivative in Assumption~\ref{ass:fixed-h} give
\begin{equation}
|R_{2n}|\lesssim_p \rho_n^2=o_p(n^{-1/2}).
\end{equation}
The same conditional argument applies uniformly on the shrinking root
neighborhood in Assumption~\ref{ass:fixed-h}. Root and numerical-Jacobian
consistency permit replacement of the intermediate derivative by $J_h$.
Rearranging yields equation~\eqref{eq:fixed-h-linearization}. The i.i.d.
central limit theorem applies to the oracle score by the $2+\delta$ moment
condition. The score difference converges in mean square, so a foldwise law of
large numbers also gives consistency of the empirical influence-function
variance. Slutsky's theorem completes the proof. This is the standard
cross-fitted Z-estimation argument in \citet{chernozhukov2018}, written here
for the two local kernel moments.

\subsection{Proof of Proposition~\ref{prop:empirical-local-rate}}

For fixed $h$, the Gaussian signals $S_h(B;\theta)$ and $K_h(B;\theta)$ are
uniformly bounded over $B$ and $\theta$. Their derivatives with respect to
$\theta$ are also uniformly bounded by constants depending only on $h$.
Consequently, the two one-dimensional translate classes indexed by the compact
set $\mathcal N_h$ have polynomial covering numbers. A maximal inequality for
bounded empirical processes therefore gives, for every training fold,
\begin{align}
\sup_{\theta\in\mathcal N_h}
|\widehat A_h^{(-k)}(\theta)-A_h(\theta)|
&=O_p\!\left(\sqrt{\frac{\log n}{n}}\right),\\
\sup_{\theta\in\mathcal N_h}
|\widehat G_h^{(-k)}(\theta)-G_h(\theta)|
&=O_p\!\left(\sqrt{\frac{\log n}{n}}\right).
\end{align}
The maximum over a fixed number of folds has the same order. These bounds can
also be obtained directly by applying Hoeffding's inequality on an
$n^{-1}$-net of $\mathcal N_h$ and using the Lipschitz envelopes between grid
points; see \citet{vandervaartwellner1996} for the general covering-number
argument. Since $\sqrt{\log n/n}=o(n^{-1/4})$, equation
\eqref{eq:empirical-local-rate} follows.

Equation~\eqref{eq:local-cdf-overlap} and uniform convergence imply
$\inf_{\theta}\widehat A_h^{(-k)}(\theta)\geq a_h/2$ for every fold with
probability approaching one. On this event, all powers of $A$ appearing in
$m_h$, $m_A$, $m_G$, and their first two nuisance derivatives are uniformly
bounded. The Gaussian signals have bounded moments for fixed $h$. Hence the
nuisance-rate and smoothness requirements stated in the proposition follow.
The remaining root conditions are the usual separated-zero conditions for a
Z-estimator; see \citet{vandervaart1998}. This proves the result.

\subsection{Proof of Proposition~\ref{prop:richardson-exact}}

The population Richardson combination satisfies
\begin{align}
r_{\mathrm R}(h)
&=\frac{\lambda^2\theta_{0,h}-\theta_{0,\lambda h}}
{\lambda^2-1}\\
&=r_0-\lambda^2b_4h^4+o(h^4),
\end{align}
because the $b_2h^2$ terms cancel. Joint asymptotic linearity at the two
bandwidths implies
\begin{equation}
\widehat r_{\mathrm R}-r_{\mathrm R}(h_n)
=\frac{1}{n}\sum_{i=1}^n
\frac{\lambda^2\varphi_{i,h_n}-\varphi_{i,\lambda h_n}}
{\lambda^2-1}
+o_p(s_{\mathrm R,n}/\sqrt n).
\end{equation}
The assumed triangular-array central limit theorem and covariance consistency
studentize this sum. Condition~\eqref{eq:richardson-rate} makes the remaining
$O(h_n^4)$ population bias negligible relative to
$s_{\mathrm R,n}/\sqrt n$. If $s_{\mathrm R,n}\asymp h_n^{-1/2}$, this
condition is $\sqrt{nh_n}\,h_n^4\to0$, equivalently $nh_n^9\to0$; the
additional condition $nh_n\to\infty$ supplies increasing local information.
